\section{Thể tích vật thể và khối tròn xoay}

\subsection{Nền tảng lý thuyết \& Định lý cơ bản}

\begin{hopdinhly}[Định lý]
Cho vật thể $(H)$ nằm giữa hai mặt phẳng vuông góc với trục $Ox$ tại $x=a$ và $x=b$. Cắt $(H)$ bởi mặt phẳng vuông góc với $Ox$ tại điểm có hoành độ $x$ ($a\le x\le b$) được thiết diện có diện tích $S(x)$, trong đó $S(x)$ liên tục trên $[a;b]$. Khi đó thể tích của vật thể $(H)$ là
\[
V = \int_a^b S(x)\dx.
\]
\end{hopdinhly}

\begin{hopchuy}[Phân biệt với công thức diện tích]
Công thức thể tích vật thể \textbf{không có $\pi$} và \textbf{không bình phương} $S(x)$ -- vì $S(x)$ \emph{đã là diện tích} của thiết diện, không phải bán kính hay hàm số $f(x)$. Đây là điểm học sinh rất hay nhầm với công thức khối tròn xoay ở Định lý~2.
\end{hopchuy}

\begin{center}
\begin{tikzpicture}[scale=1]
  \draw[navy, thick] (0,0) .. controls (1,1.6) and (3,1.8) .. (5,0.6)
        .. controls (5.3,0.3) and (5,-0.3) .. (4.6,-0.9)
        .. controls (3,-1.6) and (1,-1.3) .. (0,0) -- cycle;
  \foreach \x/\a/\b in {1.4/1.15/-1.05, 2.6/1.72/-1.15, 3.8/1.05/-0.55} {
    \draw[teal, thick] (\x,\a) ellipse [x radius=0.12, y radius=\a*0.55+0.55];
  }
  \draw[dashed, shadowgray] (1.4,-2.1) -- (1.4,1.3) node[below=2.3cm]{\small $x_1$};
  \draw[dashed, shadowgray] (2.6,-2.1) -- (2.6,2.0) node[below=2.5cm]{\small $x$};
  \draw[dashed, shadowgray] (3.8,-2.1) -- (3.8,1.3) node[below=2.3cm]{\small $x_2$};
  \draw[-Latex] (-1,-2.1) -- (6,-2.1) node[right]{$Ox$};
  \node[teal] at (2.6,2.3) {$S(x)$};
\end{tikzpicture}
\end{center}

\begin{hopdinhly}[Định lý]
Cho hàm số $y=f(x)$ liên tục, không âm trên $[a;b]$. Thể tích khối tròn xoay tạo thành khi quay hình phẳng giới hạn bởi đồ thị $y=f(x)$, trục $Ox$ và hai đường thẳng $x=a$, $x=b$ quanh trục $Ox$ là
\[
V = \pi\int_a^b f^2(x)\dx.
\]
\end{hopdinhly}

\begin{center}
\begin{tikzpicture}[scale=0.9, declare function={f(\x)=0.35*\x^2-0.9*\x+2.4;}]
  \begin{axis}[
      axis lines=middle, xlabel={$x$}, ylabel={$y$},
      xmin=-0.5, xmax=5, ymin=-1.5, ymax=3.2,
      xtick={0.5,4}, ytick=\empty, width=8.5cm, height=6cm,
      axis line style={-Latex}, samples=100, hide y axis]
    \addplot [domain=0.5:4, navy, thick, fill=navylight, fill opacity=0.35] {f(x)} \closedcycle;
    \draw[navy, thick] (axis cs:0.5,0) ellipse [x radius=0.06, y radius=1.05] ;
    \draw[navy, thick] (axis cs:4,0) ellipse [x radius=0.06, y radius=0.4] ;
  \end{axis}
\end{tikzpicture}

\small (Khối tròn xoay tạo thành khi quay miền tô màu quanh trục $Ox$)
\end{center}

\begin{hopchuy}[Vì sao không cần trị tuyệt đối]
Vì $f^2(x)\ge0$ với mọi $x$ nên $\displaystyle V=\pi\int_a^b f^2(x)\dx$ luôn không âm một cách tự nhiên -- \textbf{không cần đặt dấu trị tuyệt đối} như công thức tính diện tích.
\end{hopchuy}

\newpage
\subsection{Dạng 1: Tính thể tích vật thể có diện tích mặt cắt ngang \texorpdfstring{$S(x)$}{S(x)}}

\begin{hopphuongphap}
\begin{enumerate}[leftmargin=1cm]
  \item Dựa vào giả thiết, biểu diễn kích thước thiết diện (cạnh hình vuông, cạnh tam giác đều, bán kính nửa hình tròn,\ldots) tại hoành độ $x$ theo $x$.
  \item Dùng công thức diện tích hình học tương ứng để lập hàm số $S(x)$.
  \item Tính $\displaystyle V=\int_a^b S(x)\dx$.
\end{enumerate}
\end{hopphuongphap}

\begin{hopvidu}[(Thiết diện hình vuông)]
Cho vật thể $(H)$ có đáy là hình tròn bán kính $R$. Thiết diện của $(H)$ cắt bởi mặt phẳng vuông góc với một đường kính cố định của đáy luôn là hình vuông. Tính thể tích của $(H)$.
\end{hopvidu}
\begin{loigiai}
Chọn hệ trục $Oxy$ với gốc $O$ là tâm đường tròn đáy, trục $Ox$ trùng với đường kính cố định. Phương trình đường tròn đáy: $x^2+y^2=R^2$.

Tại hoành độ $x\in[-R;R]$, dây cung vuông góc với $Ox$ có độ dài $2y=2\sqrt{R^2-x^2}$ -- đây chính là cạnh của hình vuông thiết diện. Vậy
\[
S(x) = \left(2\sqrt{R^2-x^2}\right)^2 = 4(R^2-x^2).
\]
\[
V=\int_{-R}^{R} 4(R^2-x^2)\dx = 4\left.\left(R^2x-\frac{x^3}{3}\right)\right|_{-R}^{R} = 8\left(R^3-\frac{R^3}{3}\right).
\]
\ketqua{V=\frac{16R^3}{3}}
\end{loigiai}

\begin{hopvidu}[(Thiết diện tam giác đều)]
Cho vật thể $(H)$ có đáy là hình phẳng giới hạn bởi đường cong $y=\sqrt{x}$, trục hoành và đường thẳng $x=4$. Thiết diện của $(H)$ cắt bởi mặt phẳng vuông góc với $Ox$ tại điểm hoành độ $x$ là một tam giác đều có cạnh bằng độ dài đoạn nối từ trục $Ox$ đến đồ thị (tức bằng $y=\sqrt{x}$). Tính thể tích $(H)$.
\end{hopvidu}
\begin{loigiai}
Tam giác đều cạnh $a$ có diện tích $S=\dfrac{\sqrt3}{4}a^2$. Ở đây $a=\sqrt{x}$ nên $a^2=x$, do đó
\[
S(x)=\frac{\sqrt3}{4}x.
\]
\[
V=\int_0^4 \frac{\sqrt3}{4}x\dx = \frac{\sqrt3}{4}\left.\frac{x^2}{2}\right|_0^4 = \frac{\sqrt3}{4}\cdot8.
\]
\ketqua{V=2\sqrt3 \ (\text{đvtt})}
\end{loigiai}

\begin{hopvidu}[(Thiết diện nửa hình tròn)]
Cho vật thể $(H)$ có đáy là hình tròn bán kính $R=2$. Thiết diện của $(H)$ cắt bởi mặt phẳng vuông góc với một đường kính cố định luôn là nửa hình tròn, có đường kính bằng dây cung tương ứng của đáy. Tính thể tích $(H)$.
\end{hopvidu}
\begin{loigiai}
Với đường tròn đáy $x^2+y^2=4$, dây cung tại hoành độ $x$ có độ dài $2\sqrt{4-x^2}$, đây là \textbf{đường kính} của nửa hình tròn thiết diện, nên bán kính thiết diện là $r(x)=\sqrt{4-x^2}$. Diện tích nửa hình tròn:
\[
S(x) = \frac12\pi r^2(x) = \frac{\pi}{2}(4-x^2).
\]
\[
V=\int_{-2}^2 \frac{\pi}{2}(4-x^2)\dx = \pi\int_0^2(4-x^2)\dx = \pi\left.\left(4x-\frac{x^3}{3}\right)\right|_0^2 = \pi\left(8-\frac83\right).
\]
\ketqua{V=\frac{16\pi}{3} \ (\text{đvtt})}
\end{loigiai}

\begin{hopvidu}[(Thiết diện hình vuông từ hai đường cong)]
Cho vật thể $(H)$ có đáy là hình phẳng giới hạn bởi hai đường cong $y=\sqrt{x}$ và $y=\dfrac{x}{2}$ trên đoạn $[0;4]$. Thiết diện của $(H)$ cắt bởi mặt phẳng vuông góc với $Ox$ luôn là hình vuông có cạnh bằng bề rộng của đáy tại hoành độ đó. Tính thể tích $(H)$.
\end{hopvidu}
\begin{loigiai}
Trên $(0;4)$, ta có $\sqrt{x}\ge \dfrac{x}{2}$ (đã xét dấu ở Bài 1), nên cạnh hình vuông là
\[
a(x)=\sqrt{x}-\frac{x}{2}, \qquad S(x)=a^2(x)=x - x^{3/2}+\frac{x^2}{4}.
\]
\[
V=\int_0^4\left(x-x^{3/2}+\frac{x^2}{4}\right)\dx = \left.\left(\frac{x^2}{2}-\frac{2}{5}x^{5/2}+\frac{x^3}{12}\right)\right|_0^4 = 8-\frac{64}{5}+\frac{16}{3}.
\]
\ketqua{V=\frac{8}{15} \ (\text{đvtt})}
\end{loigiai}

\newpage
\subsection{Dạng 2: Khối tròn xoay giới hạn bởi 1 đường cong và trục \texorpdfstring{$Ox$}{Ox}}

\begin{hopphuongphap}
Áp dụng trực tiếp công thức $\displaystyle V=\pi\int_a^b f^2(x)\dx$. Kỹ năng cần rèn luyện là \textbf{bình phương $f(x)$ đúng cách} (khai triển hằng đẳng thức, dùng công thức hạ bậc với hàm lượng giác, hoặc quy tắc lũy thừa với hàm mũ) trước khi lấy nguyên hàm.
\end{hopphuongphap}

\begin{hopvidu}[(Hàm căn thức)]
Tính thể tích khối tròn xoay khi quay hình phẳng giới hạn bởi $y=\sqrt{x}$, trục $Ox$, $x=0$, $x=4$ quanh trục $Ox$.
\end{hopvidu}
\begin{loigiai}
\[
V=\pi\int_0^4 \left(\sqrt{x}\right)^2\dx = \pi\int_0^4 x\dx = \pi\left.\frac{x^2}{2}\right|_0^4.
\]
\ketqua{V=8\pi \ (\text{đvtt})}
\end{loigiai}

\begin{hopvidu}[(Hàm đa thức)]
Tính thể tích khối tròn xoay khi quay hình phẳng giới hạn bởi $y=x^2$, trục $Ox$, $x=0$, $x=2$ quanh trục $Ox$.
\end{hopvidu}
\begin{loigiai}
\[
V=\pi\int_0^2 (x^2)^2\dx = \pi\int_0^2 x^4\dx = \pi\left.\frac{x^5}{5}\right|_0^2.
\]
\ketqua{V=\frac{32\pi}{5} \ (\text{đvtt})}
\end{loigiai}

\begin{hopvidu}[(Hàm lượng giác -- công thức hạ bậc)]
Tính thể tích khối tròn xoay khi quay hình phẳng giới hạn bởi $y=\sin x$, trục $Ox$, $x=0$, $x=\pi$ quanh trục $Ox$.
\end{hopvidu}
\begin{loigiai}
Dùng công thức hạ bậc $\sin^2x = \dfrac{1-\cos2x}{2}$:
\[
V=\pi\int_0^\pi \sin^2x\dx = \pi\int_0^\pi \frac{1-\cos2x}{2}\dx = \frac{\pi}{2}\left.\left(x-\frac{\sin2x}{2}\right)\right|_0^\pi = \frac{\pi}{2}(\pi-0).
\]
\ketqua{V=\frac{\pi^2}{2} \ (\text{đvtt})}
\end{loigiai}

\begin{hopvidu}[(Hàm mũ)]
Tính thể tích khối tròn xoay khi quay hình phẳng giới hạn bởi $y=e^x$, trục $Ox$, trục $Oy$ và $x=1$ quanh trục $Ox$.
\end{hopvidu}
\begin{loigiai}
\[
V=\pi\int_0^1 (e^x)^2\dx = \pi\int_0^1 e^{2x}\dx = \pi\left.\frac{e^{2x}}{2}\right|_0^1 = \frac{\pi}{2}\left(e^2-1\right).
\]
\ketqua{V=\frac{\pi(e^2-1)}{2} \ (\text{đvtt})}
\end{loigiai}

\newpage
\subsection{Dạng 3: Khối tròn xoay tạo bởi hình phẳng giới hạn bởi 2 đường cong}

\begin{hopphuongphap}
Khi quay quanh $Ox$ hình phẳng giới hạn bởi $y=f(x)$, $y=g(x)$ (với $f,g\ge0$ và $f\ge g$ trên $[a;b]$), khối tròn xoay tạo thành có dạng \textbf{"vòng đệm"} (washer): phần thể tích do $f$ quét ra, trừ đi phần lõi do $g$ quét ra.
\[
V=\pi\int_a^b \ababs{f^2(x)-g^2(x)}\dx.
\]
\end{hopphuongphap}

\begin{hopchuy}[Sai lầm kinh điển]
Học sinh rất hay nhầm lẫn giữa hai biểu thức sau:
\[
\pi\int_a^b\ababs{f^2(x)-g^2(x)}\dx \qquad \text{(ĐÚNG)} \qquad\text{và}\qquad \pi\int_a^b (f(x)-g(x))^2\dx \qquad \text{(SAI)}.
\]
Hai biểu thức này \textbf{không bằng nhau} vì $(f-g)^2 = f^2-2fg+g^2 \ne f^2-g^2$ (trừ trường hợp $g\equiv0$). Bản chất hình học: khối tròn xoay của $(f-g)$ mô tả một vật thể hoàn toàn khác với khối vòng đệm tạo bởi $f$ và $g$.
\end{hopchuy}

\begin{hopvidu}[(Hai đường parabol/đường thẳng cơ bản)]
Tính thể tích khối tròn xoay khi quay quanh $Ox$ hình phẳng giới hạn bởi $y=x$ và $y=x^2$.
\end{hopvidu}
\begin{loigiai}
Giao điểm: $x=x^2 \Leftrightarrow x=0$ hoặc $x=1$. Trên $(0;1)$: $x\ge x^2\ge0$, nên $f(x)=x$, $g(x)=x^2$.
\[
V=\pi\int_0^1\left(x^2-x^4\right)\dx = \pi\left.\left(\frac{x^3}{3}-\frac{x^5}{5}\right)\right|_0^1 = \pi\left(\frac13-\frac15\right).
\]
\ketqua{V=\frac{2\pi}{15} \ (\text{đvtt})}
\end{loigiai}

\begin{hopvidu}[(Căn thức và đa thức)]
Tính thể tích khối tròn xoay khi quay quanh $Ox$ hình phẳng giới hạn bởi $y=\sqrt{x}$ và $y=x^2$.
\end{hopvidu}
\begin{loigiai}
Giao điểm: $\sqrt{x}=x^2 \Leftrightarrow x=0$ hoặc $x=1$. Trên $(0;1)$: $\sqrt{x}\ge x^2$ (thử $x=0{,}5$: $0{,}707>0{,}25$).
\[
V=\pi\int_0^1\left(x-x^4\right)\dx = \pi\left.\left(\frac{x^2}{2}-\frac{x^5}{5}\right)\right|_0^1 = \pi\left(\frac12-\frac15\right).
\]
\ketqua{V=\frac{3\pi}{10} \ (\text{đvtt})}
\end{loigiai}

\begin{hopvidu}[(Parabol và đường thẳng cắt nhau tại hai điểm)]
Tính thể tích khối tròn xoay khi quay quanh $Ox$ hình phẳng giới hạn bởi $(P): y=4-x^2$ và đường thẳng $d: y=x+2$.
\end{hopvidu}
\begin{loigiai}
Giao điểm: $4-x^2=x+2 \Leftrightarrow x^2+x-2=0 \Leftrightarrow x=-2$ hoặc $x=1$. Trên $(-2;1)$: tại $x=0$, $(P)=4 > d=2\ge0$, và cả hai hàm đều không âm trên đoạn này nên áp dụng trực tiếp công thức vòng đệm.
\[
f^2(x)-g^2(x) = (4-x^2)^2-(x+2)^2 = (16-8x^2+x^4)-(x^2+4x+4) = x^4-9x^2-4x+12.
\]
Gọi $F(x)=\dfrac{x^5}{5}-3x^3-2x^2+12x$. Ta có $F(1)=\dfrac{36}{5}$, $F(-2)=-\dfrac{72}{5}$, suy ra
\[
\int_{-2}^1\left(x^4-9x^2-4x+12\right)\dx = F(1)-F(-2) = \frac{36}{5}+\frac{72}{5}=\frac{108}{5}.
\]
\ketqua{V=\frac{108\pi}{5} \ (\text{đvtt})}
\end{loigiai}

\begin{hopvidu}[(Ứng dụng lại cặp hàm số đã học ở Bài 1)]
Tính thể tích khối tròn xoay khi quay quanh $Ox$ hình phẳng giới hạn bởi $y=\sqrt{x}$ và $y=\dfrac{x}{2}$ trên đoạn $[0;4]$.
\end{hopvidu}
\begin{loigiai}
Đây chính là cặp hàm số đã dùng để tính diện tích (lát gạch hồ cảnh) ở Bài 1. Trên $(0;4)$, $\sqrt{x}\ge\dfrac x2\ge0$.
\[
V=\pi\int_0^4\left(x-\frac{x^2}{4}\right)\dx = \pi\left.\left(\frac{x^2}{2}-\frac{x^3}{12}\right)\right|_0^4 = \pi\left(8-\frac{16}{3}\right).
\]
\ketqua{V=\frac{8\pi}{3} \ (\text{đvtt})}
\end{loigiai}

\newpage
\subsection{Dạng 4: Ứng dụng thực tiễn -- Tối ưu hóa trong thiết kế}

\begin{hopphuongphap}
\begin{enumerate}[leftmargin=1cm]
  \item Chọn trục quay hợp lý (thường trùng với trục đối xứng của vật thể).
  \item Xác định đường sinh (đường cong tạo ra bề mặt khi quay) dựa trên kích thước thực tế.
  \item Tính thể tích bằng công thức khối tròn xoay, sau đó đổi đơn vị nếu cần (1 dm$^3$ = 1 lít).
\end{enumerate}
\end{hopphuongphap}

\begin{hopvidu}[(Kiểm chứng công thức thể tích khối nón)]
Một hình nón có bán kính đáy $R$ và chiều cao $h$ được tạo thành bằng cách quay đoạn thẳng $y=\dfrac{R}{h}x$ ($x\in[0;h]$) quanh trục $Ox$. Dùng tích phân, hãy tính thể tích khối nón theo $R,h$ và đối chiếu với công thức đã học ở lớp dưới.
\end{hopvidu}
\begin{loigiai}
\[
V=\pi\int_0^h\left(\frac{R}{h}x\right)^2\dx = \frac{\pi R^2}{h^2}\int_0^h x^2\dx = \frac{\pi R^2}{h^2}\cdot\frac{h^3}{3}.
\]
\ketqua{V=\frac13\pi R^2 h}
Kết quả trùng khớp với công thức thể tích khối nón đã học ở lớp dưới -- đây là một minh chứng cho tính đúng đắn của phương pháp tích phân.
\end{loigiai}

\begin{hopvidu}[(Thùng đựng rượu vang)]
Một thùng đựng rượu có dạng khối tròn xoay, tạo thành khi quay hình phẳng giới hạn bởi parabol $y=2-0{,}5x^2$ (đơn vị dm) và trục hoành quanh trục $Ox$, với $x\in[-2;2]$. Tính thể tích thùng (theo lít, biết $1\text{dm}^3=1$ lít).
\end{hopvidu}
\begin{loigiai}
\[
y^2=(2-0{,}5x^2)^2 = 4-2x^2+0{,}25x^4.
\]
\begin{align*}
V&=\pi\int_{-2}^{2}\left(4-2x^2+0{,}25x^4\right)\dx = 2\pi\int_0^2\left(4-2x^2+0{,}25x^4\right)\dx \\
 &= 2\pi\left.\left(4x-\frac{2x^3}{3}+0{,}05x^5\right)\right|_0^2 = 2\pi\left(8-\frac{16}{3}+1{,}6\right).
\end{align*}
\ketqua{V=\frac{128\pi}{15}\approx26{,}8 \text{ lít}}
\end{loigiai}

\begin{hopvidu}[(Bồn chứa hình chỏm parabol)]
Một chiếc chậu có dạng khối tròn xoay, tạo thành khi quay đường cong $x=\dfrac{\sqrt{y}}{2}$ (đơn vị dm) quanh trục $Oy$, với $0\le y\le 4$. Tính thể tích của chậu (theo lít).
\end{hopvidu}
\begin{loigiai}
Khi quay quanh trục $Oy$, công thức thể tích tương tự Định lý 2 nhưng đổi vai trò $x\leftrightarrow y$:
\[
V=\pi\int_0^4 x^2\dif{y} = \pi\int_0^4 \frac{y}{4}\dif{y} = \frac{\pi}{4}\left.\frac{y^2}{2}\right|_0^4 = \frac{\pi}{4}\cdot8.
\]
\ketqua{V=2\pi\approx6{,}28 \text{ lít}}
\end{loigiai}

\begin{hopvidu}[(Ly thủy tinh hình cong)]
Một chiếc ly có thành trong dạng khối tròn xoay, tạo thành khi quay đường cong $x=\dfrac{y^2}{16}$ (đơn vị cm) quanh trục $Oy$, với $0\le y\le8$ ($y$ là chiều cao tính từ đáy ly). Tính thể tích ly (theo ml, biết $1\text{cm}^3=1$ ml).
\end{hopvidu}
\begin{loigiai}
\[
V=\pi\int_0^8 x^2\dif{y} = \pi\int_0^8 \frac{y^4}{256}\dif{y} = \frac{\pi}{256}\left.\frac{y^5}{5}\right|_0^8 = \frac{\pi}{256}\cdot\frac{32768}{5}.
\]
\ketqua{V=\frac{128\pi}{5}\approx80{,}4 \text{ ml}}
\end{loigiai}

\newpage
\subsection{Bài tập tự luyện}

\subsubsection*{Phần I -- Trắc nghiệm nhiều phương án lựa chọn}

\begin{multicols}{2}
\cauhoi Thể tích khối tròn xoay khi quay hình phẳng giới hạn bởi $y=x$, $Ox$, $x=0$, $x=2$ quanh $Ox$ bằng
\begin{enumerate}[label=\textbf{\Alph*.}, leftmargin=0.9cm]
  \item $\dfrac{8\pi}{3}$
  \item $\dfrac{4\pi}{3}$
  \item $2\pi$
  \item $\dfrac{8}{3}$
\end{enumerate}

\cauhoi Thể tích vật thể có $S(x)=x^2+1$ trên $[0;2]$ bằng
\begin{enumerate}[label=\textbf{\Alph*.}, leftmargin=0.9cm]
  \item $\dfrac{14}{3}$
  \item $\dfrac{8}{3}$
  \item $6$
  \item $\dfrac{10}{3}$
\end{enumerate}

\cauhoi Quay hình phẳng giới hạn bởi $y=\cos x$, $Ox$, $x=0$, $x=\dfrac{\pi}{2}$ quanh $Ox$, thể tích khối tròn xoay bằng
\begin{enumerate}[label=\textbf{\Alph*.}, leftmargin=0.9cm]
  \item $\dfrac{\pi^2}{4}$
  \item $\dfrac{\pi}{2}$
  \item $\pi^2$
  \item $\dfrac{\pi^2}{2}$
\end{enumerate}

\cauhoi Thể tích khối cầu bán kính $R$ (tạo bởi quay nửa đường tròn $y=\sqrt{R^2-x^2}$ quanh $Ox$) bằng
\begin{enumerate}[label=\textbf{\Alph*.}, leftmargin=0.9cm]
  \item $\dfrac{4}{3}\pi R^3$
  \item $\dfrac{2}{3}\pi R^3$
  \item $4\pi R^2$
  \item $\pi R^3$
\end{enumerate}

\cauhoi Quay hình phẳng giới hạn bởi $y=e^x$, $Ox$, $Oy$, $x=1$ quanh $Ox$, thể tích bằng
\begin{enumerate}[label=\textbf{\Alph*.}, leftmargin=0.9cm]
  \item $\dfrac{\pi(e^2-1)}{2}$
  \item $\pi(e-1)$
  \item $\dfrac{\pi e^2}{2}$
  \item $\pi(e^2-1)$
\end{enumerate}

\cauhoi Công thức đúng để tính thể tích khối tròn xoay giới hạn bởi 2 đường cong $f(x)\ge g(x)\ge0$ trên $[a;b]$, quay quanh $Ox$, là
\begin{enumerate}[label=\textbf{\Alph*.}, leftmargin=0.9cm]
  \item $\pi\displaystyle\int_a^b[f(x)-g(x)]^2\dx$
  \item $\pi\displaystyle\int_a^b[f^2(x)-g^2(x)]\dx$
  \item $\displaystyle\int_a^b[f^2(x)-g^2(x)]\dx$
  \item $\pi\displaystyle\int_a^b[f(x)-g(x)]\dx$
\end{enumerate}

\cauhoi Thể tích khối tròn xoay giới hạn bởi $y=x^2$ và $y=4$, quay quanh $Ox$ (trên miền $y=x^2$ nằm dưới $y=4$) bằng
\begin{enumerate}[label=\textbf{\Alph*.}, leftmargin=0.9cm]
  \item $\dfrac{512\pi}{5}$
  \item $\dfrac{256\pi}{5}$
  \item $\dfrac{128\pi}{5}$
  \item $32\pi$
\end{enumerate}

\cauhoi Vật thể có đáy là hình tròn bán kính $2$, thiết diện vuông góc một đường kính cố định luôn là tam giác đều. Thể tích vật thể bằng
\begin{enumerate}[label=\textbf{\Alph*.}, leftmargin=0.9cm]
  \item $\dfrac{32\sqrt3}{3}$
  \item $\dfrac{16\sqrt3}{3}$
  \item $8\sqrt3$
  \item $\dfrac{64\sqrt3}{3}$
\end{enumerate}

\cauhoi Quay hình phẳng giới hạn bởi $y=\sqrt{x}$, $y=x$ quanh $Ox$, thể tích khối tròn xoay bằng
\begin{enumerate}[label=\textbf{\Alph*.}, leftmargin=0.9cm]
  \item $\dfrac{\pi}{6}$
  \item $\dfrac{\pi}{3}$
  \item $\dfrac{2\pi}{15}$
  \item $\dfrac{\pi}{2}$
\end{enumerate}

\cauhoi Một hình nón có bán kính đáy $3$ và chiều cao $4$. Thể tích khối nón (tính bằng tích phân) bằng
\begin{enumerate}[label=\textbf{\Alph*.}, leftmargin=0.9cm]
  \item $12\pi$
  \item $36\pi$
  \item $16\pi$
  \item $4\pi$
\end{enumerate}
\end{multicols}

\subsubsection*{Phần II -- Trắc nghiệm Đúng/Sai}

\cauhoi Cho hình phẳng $(D)$ giới hạn bởi $y=x^2$ và $y=2x$.
\begin{enumerate}[label=\textbf{\alph*)}, leftmargin=1cm]
  \item Hoành độ giao điểm của hai đồ thị là $x=0$ và $x=2$.
  \item Trên $(0;2)$, $2x\ge x^2$.
  \item Thể tích khối tròn xoay khi quay $(D)$ quanh $Ox$ được tính bởi $\pi\displaystyle\int_0^2(2x-x^2)^2\dx$.
  \item Thể tích khối tròn xoay khi quay $(D)$ quanh $Ox$ bằng $\dfrac{64\pi}{15}$.
\end{enumerate}

\cauhoi Cho vật thể $(H)$ có đáy là hình tròn bán kính $1$, thiết diện vuông góc một đường kính cố định luôn là hình vuông.
\begin{enumerate}[label=\textbf{\alph*)}, leftmargin=1cm]
  \item Diện tích thiết diện tại hoành độ $x$ là $S(x)=4(1-x^2)$.
  \item Thể tích vật thể được tính bởi $\pi\displaystyle\int_{-1}^1 S(x)\dx$.
  \item Thể tích vật thể bằng $\dfrac{16}{3}$.
  \item Nếu thay thiết diện hình vuông bằng nửa hình tròn thì thể tích sẽ giảm đi.
\end{enumerate}

\subsubsection*{Phần III -- Trắc nghiệm trả lời ngắn}

\cauhoi Tính thể tích khối tròn xoay khi quay hình phẳng giới hạn bởi $y=\ln x$, $Ox$, $x=e$ quanh trục $Ox$ (kết quả theo $\pi$, làm tròn hệ số đến hàng phần trăm nếu cần).

\cauhoi Một chiếc phễu có dạng hình nón (tạo bởi quay đoạn thẳng $y=2x$, $x\in[0;5]$ quanh $Ox$, đơn vị cm). Tính thể tích phễu (theo cm$^3$, làm tròn đến hàng đơn vị, lấy $\pi\approx3{,}14$).

\cauhoi Vật thể có đáy là hình phẳng giới hạn bởi $y=4-x^2$ và trục hoành, thiết diện vuông góc trục $Ox$ luôn là hình vuông có cạnh bằng $\dfrac12$ bề rộng đáy tại hoành độ đó. Tính thể tích vật thể (kết quả dạng phân số tối giản).

\cauhoi Một bể cá hình trụ tròn xoay được tạo bởi quay hình chữ nhật có các cạnh $y=3$ (dm), $x\in[0;5]$ quanh trục $Ox$. Tính thể tích bể (theo lít).

\clearpage
\begin{hopchuy}[Đáp án]
\textbf{Phần I:} 1-A, 2-A, 3-A, 4-A, 5-A, 6-B, 7-B, 8-A, 9-A, 10-A.

\textbf{Phần II:} Câu 11: a-Đ, b-Đ, c-S (phải là $\pi\displaystyle\int_0^2(4x^2-x^4)\dx$, không phải $(2x-x^2)^2$), d-Đ (vì $\pi\displaystyle\int_0^2(4x^2-x^4)\dx=\pi\left(\dfrac{32}{3}-\dfrac{32}{5}\right)=\dfrac{64\pi}{15}$). \quad Câu 12: a-Đ, b-S (công thức $V=\displaystyle\int_{-1}^1 S(x)\dx$ không nhân $\pi$), c-Đ, d-Đ (nửa hình tròn cho $V=\dfrac{2\pi}{3}\approx2{,}09 < \dfrac{16}{3}\approx5{,}33$).

\textbf{Phần III:} Câu 13: $V=\pi(e-2)$ (đvtt) $\approx0{,}72\pi$. \quad Câu 14: $V=\dfrac13\pi(10)^2(5)=\dfrac{500\pi}{3}\approx523{,}6\Rightarrow524$ cm$^3$. \quad Câu 15: $V=\dfrac{128}{15}$. \quad Câu 16: $V=\pi\cdot3^2\cdot5=45\pi\approx141{,}4$ lít.
\end{hopchuy}

\newpage
